\begin{subsection}{Put a ring on it: Varieties and their local study}
    In this section we define the basic tools of algebraic geometry, such as 
    vanishing sets/ideals, varieties and the structure sheaf of a variety. 
    All of these constructions have an analoge in the setting of risk measures as we will see in the latter chapters.
    We won't cover the information, which can be extracted from the structure sheaf and its constitutents, such as the dimension
    of the variety, tangent space and singularities, since we won't be able to do this for the risk measure case, but the interested
    reader is referred to Milne's excellent lecture \cite{milne2012algebraic} or Bosch's \cite{alggeoandcom}.

    %Thus in the subsequent section we will first speak in some generality of subrings 
    Let \( R \) be a commutative ring resulting from taking a quotient of \( \stdpolyr \) with respect to some ideal \( I \).
    \begin{definition}
        For a subset \( J \subset \stdpolyr \) and \( W \subset \stdaff \) define
        \begin{align*}
            V(J) &:= \left\{ x \in \stdaff \mid \forall f \in J,\ f(x) = 0 \right\}  \\
            I(W) &:= \left\{ f \in \stdpolyr \mid \forall x \in W,\ f(x) = 0 \right\}
        \end{align*}
        as the vanishing set and vanishing ideal respectively. Complementarily, let
        \( D(J) = \stdaff \setminus V(J) \) be the non-vanishing set. In the case of principal ideals \( J = (f) \)
        we call \( D(f) := D((f)) \) a distinguished open set and it holds \( D(fg) = D(f) \cap D(g) \).
    \end{definition}
    That \( I(W) \) is indeed an ideal in \( \stdpolyr \) is a standard verification.
    \begin{proposition}
        The operation \( V \) is inclusion reversing, meaning for \( J_1 \subseteq J_2 \)
        it holds \( V(J_2) \subseteq V(J_1) \) . 
        Further for any collection \( \{J_n\}_{n \in N} \) with \( J_n \subset \stdpolyr \)
        \begin{align*}
            V\left(\bigcup_{n \in N} J_n \right) = \bigcap_{n \in N} V\left(J_n \right),
        \end{align*}
        but in general only
        \begin{align*}
            V\left(\bigcap_{n \in N} J_n \right) \subseteq \bigcup_{n \in N} V\left(J_n \right).
        \end{align*}
        The same is true of \( I \) with sets in \( \stdaff \). The dual statements hold for the non-vanishing sets \( D(J) \).
    \end{proposition}
    \begin{proof}
        Proposition 6.1.1 of \cite{alggeoandcom}.
    \end{proof}
    
    \begin{definition}
        An algebraic set \( W \subseteq \stdaff \) is a set of the form \( W = V(J) \) for some \( J \subset \stdpolyr \).
        A variety \( W \) is an irreducible algebraic set, meaning there are no algebraic sets \( W_1, W_2 \), s.t.
        \( W = W_1 \cup W_2 \).
    \end{definition}
    \begin{proposition} \label{primenessandradicalag}
        We have \( W \subseteq \stdaff \) is a variety if and only if \( I(W) \) is a prime ideal or the whole ring \( \stdpolyr \).
        Further for any ideal \( J \subseteq \stdpolyr \) it holds \( V(\sqrt{J}) = V(J) \).
    \end{proposition}
    \begin{proof}
        Corollary 6.1.7 of \cite{alggeoandcom} and Corollary 6.1.16 of \cite{alggeoandcom}.
    \end{proof}
    
    \begin{example}
        Consider \( \aff{\mathbb{R}}{2} \) then the unit circle is an algebraic variety induced by the principal ideal \( (x_1^2 + x_2^2 - 1) \),
        while the cross of the coordinate axes, denoted by \( W \), is not since \( W = V(\{x_1\}) \cup V((x_2)) \). Note also that the affine space
        itself is a variety, since it is induced by the ideal \( (0) \), which is naturally prime and any point \( c \in \stdaff \) via 
        \( \{c\} = V((x_1 - c_1), \dots (x_n - c_n)) \). The fact that points are themselves varieties leads to the more general study
        of algebraic geometry, where the prime ideals are taken as the points of the "geometric space" spectrum \( \spec{\stdaff} \).
        We won't do this abstraction as we will later see in the case of risk measure, that this gives much more points then we're interested in.
    \end{example}
    \begin{definition}
        The Zariski topology on \( \stdaff \) is the topology induced by the distinguished open sets \( D(f) \)
        as a subbasis of open sets.
    \end{definition}
    \newline
    Now an implicit abstraction that might be confusing should be pointed out. Studying the local behavior
    of some algebraic structure around a geometric object is always a relative matter. Meaning
    there are two different positions involved, first the one of the geometric object and second the one were the local study
    is conducted. In the simplest case the geometric space is the whole space thus the local study is unperturbed by any effects 
    propagated by any special geometry of its overall surroundings. Thus we have one algebraic structure of the whole space which needs 
    to be adjusted when focusing on geometric object. However, in both cases the procedure of localizing it should be the same mechanism.
    More concretely, we have to differentiate between the local study on the whole affine space \( \stdaff \)
    vs. the study restricted on some variety \( W \subset \stdaff \). On the algebraic 
    side this is reflected if we \textit{localize} on the whole multivariate polynomial ring \( \stdpolyr \) or the quotient ring
    \( \mathbb{K}[W] := \stdpolyr / I(W) \) of it. Taking the quotient ring is natural, because all the polynomials which vanish over \( W \) can not add  
    any further information, so we would like to concentrate on those elements of the ring unaffected by them. Note also, that there is no issue
    in evaluating elements \( g_1, g_2 \) in the same equivalence class \( [g] \in \mathbb{K}[W] \) on points in \( W \), since by assumption there
    is some \( h \in I(W) \) s.t. \( g_1 + h = g_2 \) and for any point \( x \in W \), \( g_1(x) + h(x) = g_1(x) + 0 = g_2(x) \).
    Further \( W \) inherits naturally the subspace topology of \( \stdaff \) and one can verify this is the same as inducing it by the subbasis \( D([g]) \)
    for \( [g] \in \mathbb{K}[W] \). Thus we can introduce relative notions of the operations \( V, I, D \) as follows 
    \begin{align*}
        V_{W}(J) &:= \left\{ x \in W \mid \forall [f] \in J,\ [f(x)] = 0 \right\}  \\
        I_{W}(U) &:= \left\{ [f] \in \mathbb{K}[W] \mid \forall x \in U,\ [f(x)] = 0 \right\},
    \end{align*}
    where \( J \subset \mathbb{K}[W] \), \( U \subset W \) and \( D_W(J) \) as the complement of \( V_{W}(J) \) in \( W \).
    \begin{definition}
        Let \( W \subset \stdaff \) be a variety and  \( S \subset \mathbb{K}[W] \) a set that is multiplicatively closed meaning \( S \cdot S = S \).
        %Let \( S_W = \stdpolyr \setminus I(\{ W \}) \) 
        The localization at \( U \subseteq W \) is defined as the ring, denoted by \( \mathbb{K}[W]_{U} \), resulting from the localization of \( \mathbb{K}[W] \) 
        with the set \( S_U := \mathbb{K}[W] \setminus I_{W}(U) \), if it is multiplicative.
    \end{definition}
    One can make a set \( S \) multiplicatively closed by taking its saturation, which means adding all elements from the ring s.t. their product is in \( S \).
    A standard fact of algebra is that the prime ideals in \( \mathbb{K}[W]_{U} \) are all prime ideals in \( \mathbb{K}[W] \) which are contained in \( I_W(U) \),
    thus by Proposition \ref{primenessandradicalag} the concept of subvarieties inside \( W \) is well-defined algebraically in \( \mathbb{K}[W]_{U} \).
    %Note that in the special case of \( U = D_W([f]) \) this is precisely the ring \( S^{-1}_{U} \mathbb{K}[W] \) where all elements of the form \( [f^m] \)
    %are invertible.
    \begin{theorem}[Structure Sheaf] \label{sheafofvariety}
        % TODO is this well-defined since we are localizing at general U (whose ideals are not necessarily complements of prime ideals)
        Let \( W \subset \stdaff \) be a variety with its inherited Zariski topology, the map 
        \begin{align*}
            \mathcal{O}_W : \text{Open}(W) & \to \text{Ring} \\
            U & \mapsto \varinjlim_{U' \in \mathcal{I}} \mathbb{K}[W]_{D_W([f])}
        \end{align*}
        defines a sheaf of rings on \( W \), where \( \mathcal{I} \) is the set of distinguished open sets \( D_W([f]) \) s.t. \( D_W([f]) \subset U \). 
        Sections of \( \mathcal{O}(U) \) are referred to as regular functions. 
    \end{theorem}
    \begin{proof}
        A special case of Theorem 6.6.3 and Lemma 6.6.4 in \cite{alggeoandcom}.
    \end{proof}
    Explicitely, one can describe \( \mathcal{O}_W(U) \) as the set of functions \( \varphi \) s.t. for some \( x \in U \) there is some open neighborhood \( U' \subseteq U \)
    s.t. \( \varphi|_{U'} = \frac{f}{g} \) where \( f, g \in \stdpolyr \) and \( g \notin V(U') \). 
    A non-trivial intermediate in the proof of Theorem \ref{sheafofvariety} is that for \( U = D_W([f]) \) we have \( \mathbb{K}[W]_{D_W([f)]} = (f)^{-1} \mathbb{K}[W] \),
    meaning that the sections are simply fractions \( \left[\frac{h}{f^k}\right] \), fitting the previous description.
    Further for distinguished open sets \( D_W([g]) \subset D_W([f]) \) one can derive from Proposition \ref{primenessandradicalag} 
    that \( [g^m] = [h \cdot f] \) for some \( h \in \stdpolyr \) and therefore any \( \left[\frac{b}{f^k}\right] \in \mathcal{O}_W(D_W([f])) \) can be written as 
    \( \frac{b h^k}{g^{mk}} \) and the restriction morphism \( r_{D_W([g]),D_W([f])} \) simply rewrites fractions with \( f \) in the denominator into ones with \( g \).
    Taking the colimit corresponds roughly to writing all fractions in the individual distuingished sets in terms of their common denominators, which do not vanish over \( U \),
    resulting in the ring of regular functions. The matter that the \( \varphi \) are only locally of this form is precisely the side-effect of taking this limit.
    Moreover, the stalk \( \mathcal{O}_{W,x} \) are all rational polynomials \( \frac{f}{g} \) s.t. \( g(x) \neq 0 \), the stalk is in fact a local ring,
    whose unique maximal ideal \( \mathfrak{m}_x \) is given by all polynomials which vanish at \( x \).
    Effectively, these rational polynomials are only defined with respect to a point we could identify them with their value there and this 
    is the idea of the residue field.
    %In the case of algebraic geometry there is one further technique possible for the sheaf of Theorem \ref{sheafofvariety}, namely the residue field,
    %which are the aforementioned evaluations.
    \begin{definition}
        The residue field at \( x \in W \) is defined as the quotient field obtained from \( \kappa(x) = \mathcal{O}_W(x) / \mathfrak{m}_x  \).
    \end{definition}
    Since taking the quotient field with respect to the ideal of a point \( (x_1 - c_1, \dots, x_n - c_n) \), replaces each variable in a polynomial
    by the corresponding \( c_n \), we obtain the value of \( \frac{f}{g} \) at \( x \).

    \iffalse

    In the case of the variety \( W = \stdaff \), the whole space, the quotient ring coincides with \( \stdpolyr \) thus we do not have to be careful about the  
    equivalence classes. In fact this is also not much of an issue in the case of a variety \( W \) as one can see below 
    \begin{corollary}
        For a variety \( W \subset \stdaff \) and any open \( U \subset \stdaff \) the map 
        \[ \mathcal{I}_W(U) = \left\{f \in \mathcal{O}_{\stdaff}(U) \mid \forall x \in U \cap W, [f(x)] = 0 \right\} \] 
        for some open \( U \subset \stdaff \) in the Zariski topology, defines a sheaf of rings as well.
        Let \( W \subset \stdaff \) be a variety, then for any \( U \subseteq \stdaff \) there is a short exact sequence 
        \begin{align}
            0 \to \mathcal{I}_{\stdaff}(U) \to \mathcal{O}_{\stdaff}(U) \to i_\ast \mathcal{O}_W(U \cap W) \to 0
        \end{align}
        where \( i_\ast \) is the pushforward of the sheaf \( \mathcal{O}_W \) via the inclusion map \( i : W \to \stdaff \), 
        in other words \( \mathcal{O}_W(U) \simeq \mathcal{O}_{\stdaff}(U) / \mathcal{I}_{\stdaff}(U) \).
    \end{corollary}
    \begin{proof}
        By Definition \( \mathcal{I}_{\stdaff}(U) \) is a subring of \( \mathcal{O}_{\stdaff}(U) \), thus the injective part of the sequence is clear.
        Since any open \(  \)
        The morphism \( \mathcal{O}_{\stdaff}(U) \to i_\ast \mathcal{O}_W(U) \) is defined on distuingished opens \( D(f) \) 
        For the surjective one 
    \end{proof}

    The previous observation is made precise by the following

    \begin{corollary}
        Let \( W \subset \stdaff \) be a variety, then there is a short exact sequence 
        \begin{align}
            0 \to \mathcal{I}_W \to \mathcal{O}_{\stdaff} \to i_\ast \mathcal{O}_W \to 0
        \end{align}
        where \( i_\ast \) is the pushforward of the sheaf \( \mathcal{O}_W \) via the inclusion map \( i : W \to \stdaff \)
    \end{corollary}

    \begin{definition}
        Let \( W \susbet \stdaff \) be a variety, then we call the localization at \( \mathcal{O}_{\stdaff}(W) \) the coordinate ring of \( W \), denoted by \( \mathcal{O}_{\stdaff} \).
        be the multiplicative closed set induced by a set \( W \subseteq \stdaff \).
        The localization \( \mathcal{O}_{\stdaff}(W) \) at \( W \subseteq \stdaff \) is defined as the ring resulting from the localization of the ring \( \stdpolyr \) 
        with the set \( S \). 
    \end{definition}
    

    \begin{theorem}

    \end{theorem}
    One can then show 
    \begin{definition}
        For a subset \( J \subset \stdpolyr \) and \( W \subset \stdaff \) define
        \begin{align*}
            V(J) := \left\{ \mathfrak{p} \in \spec{\stdpolyr} \mid \forall f \in J, f \in \spec{\stdpolyr} \right\}  \\
            I(W) := \left\{ f \in \stdpolyr \mid \forall x \in W,\ f(x) = 0 \right\}
        \end{align*}
        as the vanishing set and vanishing ideal respectively. Complementarily, let
        \( D(J) = \stdaff \setminus V(J) \) be the non-vanishing set.
    \end{definition}
    \begin{definition}
        Call \( D(f) \) for fixed \( f \in \stdpolyr \) an elementary non-vanishing set.
        The Zariski topology on \( \spec{\stdpolyr} \) is the topology induced by the \( D(f) \)
        as a subbasis of open sets.
    \end{definition}
    \iffalse
    In particular for the one can deduce from 
    the relative version of 
    \( D([g]) \subset D([f]) \) it holds  that 

    \( \mathcal{O}_W(U) \) are also rings, since adding or multiplying two non-vanishing functions \( f_1, f_2 \) yields another one over the intersection \( U_1 \cap U_2 \) of their domains, 
    turning \( \aff{n}{\mathbb{K}} \) into a locally ringed space. In fact this is the inspiration of schemes, which are sheaves of rings that are isomorphic to the previous sheaf  
    construction but for more general rings then \( \stdpolyr \) or its quotient rings.
    %One could define schemes in the sense of Grothendieck as exactly these objects (which is far less general than the usual definition into the category of rings, which is isomorphic  
    \fi
    \fi

\end{subsection}
